% Compilar com: latexmk -lualatex entrega.tex  (depois de compilar artigo/artigo.tex)
\documentclass[11pt,a4paper]{article}
\usepackage[brazil]{babel}
\usepackage{fontspec}
% Fallback para símbolos que o Latin Modern não tem (⌈ ⌉ ≤ → ✓ ...), comuns nos logs
\directlua{luaotfload.add_fallback("simbolos", {"DejaVu Sans:mode=harf;", "DejaVu Sans Mono:mode=harf;", "FreeSerif:mode=harf;"})}
\setmainfont{Latin Modern Roman}[RawFeature={fallback=simbolos}]
\setmonofont{DejaVu Sans Mono}[Scale=0.85]
\usepackage[margin=2.2cm]{geometry}
\usepackage{pdfpages}
\usepackage{xcolor}
\usepackage[hidelinks]{hyperref}
\usepackage[pipeTables, tableCaptions, fencedCode, smartEllipses, hashEnumerators]{markdown}
\usepackage{array}
\usepackage{xurl}
\usepackage{fvextra}  % quebra de linha nos blocos de código do log
\fvset{breaklines=true, breakanywhere=true, fontsize=\small}
% "####" vira título em bloco (o \paragraph padrão fica na mesma linha e cola na tabela seguinte)
\markdownSetup{renderers={headingFour={\par\medskip\noindent\textbf{#1}\par\nopagebreak}}}
% Tabelas do log: colunas de largura igual com quebra de linha, em vez de "l" (que estoura a margem)
\makeatletter
\def\markdownLaTeXReadAlignments#1{%
  \advance\markdownLaTeXColumnCounter by 1\relax
  \addto@hook\markdownLaTeXTableAlignment{>{\raggedright\arraybackslash\footnotesize}p{\dimexpr(\linewidth-2\tabcolsep*\markdownLaTeXColumnTotal)/\markdownLaTeXColumnTotal\relax}}%
  \ifnum\markdownLaTeXColumnCounter<\markdownLaTeXColumnTotal\relax\else
    \expandafter\@gobble
  \fi\markdownLaTeXReadAlignments}
\makeatother
\newcommand{\revisar}[1]{\textcolor{red}{[REVISAR: #1]}}
\setcounter{secnumdepth}{0}
\setlength{\parindent}{0pt}
\setlength{\parskip}{0.5em}
\sloppy

\begin{document}

% ---------------- Capa (mesmo modelo do EP1) ----------------
\begin{titlepage}
\centering
Universidade de São Paulo\\
Escola de Artes, Ciências e Humanidades (EACH)
\vspace*{\fill}

{\LARGE\bfseries Exercício Prático 2\par}
\vspace{0.6em}
{\large Redes de Computadores\par}
\vspace{0.6em}
{\large Revisão crítica de literatura: lookup de rota por\\ \textit{longest prefix match}, TCAM e algoritmos\par}
\vspace*{\fill}

Gabriel Trigueiro Melo -- 15511832\\
Ana Vitória Chaves Viana -- 12733869\\
Diogo dos Santos da Rocha -- 13673281
\vspace*{\fill}

São Paulo\\
\revisar{data da entrega}
\end{titlepage}

% ---------------- Parte I ----------------
\includepdf[pages=-,pagecommand={\thispagestyle{empty}}]{../artigo/artigo.pdf}

% ---------------- Parte II ----------------
\clearpage
\part*{Parte II: Log completo de conversas com a IA}
\addcontentsline{toc}{part}{Parte II: Log completo de conversas com a IA}

Este log reúne todas as conversas do grupo com a IA durante o trabalho, na ordem em que aconteceram.
\begin{itemize}
\item \textbf{31/10 e 01/11/2026, claude.ai:} os dois primeiros arquivos foram copiados sem alteração dos registros
  feitos por Gabriel no Projeto ``EP 2 - Redes''. Lá as ações da IA (buscas, leitura de arquivos) aparecem
  como linhas em itálico marcadas \emph{Ações da IA}.
\item \textbf{A partir de 16/11/2026, Claude Code:} cada sessão foi exportada automaticamente do transcript
  que o Claude Code grava (arquivo JSONL), pelo script \texttt{.claude/skills/log-ia/export\_log.py} do
  repositório. Os prompts e as respostas em texto aparecem inteiros. As chamadas de ferramentas aparecem
  resumidas em itálico, como \emph{Ações da IA}. Para cada subagente aparecem o pedido enviado e o relatório
  devolvido. Quando o transcript guardou uma fala só como resumo de uma linha, ela aparece marcada como
  \emph{(narração)}.
\item \textbf{O que não aparece:} as instruções internas do sistema, o raciocínio interno do modelo e a
  saída bruta das ferramentas (conteúdo de arquivos e comandos).
\end{itemize}

\markdownInput{../logs/01_claude_ai_31_10_2026.md}
\clearpage
\markdownInput{../logs/02_correcoes_01_11_2026.md}
\clearpage
\markdownInput{../logs/cc_20261116_1050_gabriel.md}
\clearpage


% ---------------- Parte III ----------------
\clearpage
\part*{Parte III: Declaração de uso de IA}
\revisar{declaração a ser escrita na Fase 7}

\end{document}
